Tabular data remains one of the most important data formats in practical machine learning. Many high-value prediction problems, such as risk assessment and health-related prediction, are naturally represented as tables, where rows are observations and columns are structured variables. Yet supervised learning on tabular data has progressed more slowly than in vision or language, and strong tree-based methods such as GBDTs remain difficult to displace in practice~\cite{ref1}.

A key reason is that tabular data has a different structure from images and text. In images, nearby pixels have local spatial relationships~\cite{ref2}, and in language, tokens derive meaning from sequence context~\cite{ref3}. By contrast, tabular features often have standalone semantic meaning, and different parts of the input space may depend on different predictive variables~\cite{ref1}. This makes it hard for a single global model to capture all relevant patterns well, especially on heterogeneous datasets~\cite{ref1}.

xRFM is designed to address this gap by combining a feature-learning kernel model with an adaptive binary tree~\cite{ref1}. At each internal node, it uses the Average Gradient Outer Product (AGOP) to choose a split direction, and at each leaf it fits a local Recursive Feature Machine~\cite{ref1,ref4}. This is interesting for two reasons: it allows feature relevance to vary across regions of the data space, and it improves scalability by replacing one large global kernel problem with a hierarchy of smaller local ones~\cite{ref1}.

In this project, we evaluate xRFM on five tabular datasets---Abalone, Diamonds, Superconductor, Obesity Level, and Sleep Health---covering three regression tasks, two classification tasks, mixed-type features, a high-dimensional dataset, and large-sample settings. We compare it with XGBoost, Random Forest, and TabPFN, and also study AGOP-based interpretability and learning-curve behaviour. Overall, our results show that xRFM is a competitive and interpretable method: it performs strongly on several tasks and is often close to the best baseline, but it does not consistently outperform TabPFN or XGBoost across all datasets, especially on the more challenging large-scale classification setting.